\subsection{Symmetric tensors over a free module}

PROPOSITION 4. --- Let M be free and \((e_{i})_{i\in I}\) a basis of M.

\begin{enumerate}
\def\labelenumi{(\roman{enumi})}
\tightlist
\item
  For \(\nu \in \mathbf{N}^{(I)}\) let \(e_{\nu} = \prod_{i\in I}\gamma_{\nu_i}(e_i)\) . Then \((e_{\nu})_{\nu \in \mathbf{N}^{(I)}}\) is a basis of the A-module
\end{enumerate}

\(\mathbf{TS}(\mathbf{M})\) . In particular the algebra \(\mathbf{TS}(\mathbf{M})\) is generated by the family of elements \(\gamma_k(x)\) for \(k \in \mathbf{N}\) and \(x \in \mathbf{M}\) .

\begin{enumerate}
\def\labelenumi{(\roman{enumi})}
\setcounter{enumi}{1}
\tightlist
\item
  For each \(p \in \mathbf{N}\) , \(\mathbf{T}\mathbf{S}^p(\mathbf{M})\) is a direct factor of the A-module \(\mathbf{T}^p(\mathbf{M})\) .
\end{enumerate}

Let us use the notation of the Remark 2 above. The family \((e_{\rho(1)} \otimes \ldots \otimes e_{\rho(p)})_{\rho \in \mathcal{M}}\) is a basis of \(\mathbf{T}^p(\mathbf{M})\) . Hence Prop. 4 follows from formula (7) and the following lemma, applied with \(\mathbf{H} = \mathfrak{S}_p\) and \(\mathbf{U} = \mathbf{T}^p(\mathbf{M})\) .

Lemma 1. --- Let H be a finite group and U a left A{[}H{]}-module. Suppose that the A-module U has a basis B which is stable under the operations of H in U, and put \(\Omega = \mathrm{B} / \mathrm{H}\) . For each \(\omega \in \Omega\) let \(u_{\omega} = \sum b\) ; then

\begin{enumerate}
\def\labelenumi{(\roman{enumi})}
\item
  \((u_{\omega})_{\omega \in \Omega}\) is a basis of the A-module \(\mathbf{U}^{\mathrm{H}}\) .
\item
  For each \(\omega \in \Omega\) let \(v_{\omega}\) be a point of \(\omega\) ; put \(\omega' = \omega - \{v_{\omega}\}\) and \(\mathbf{B}' = \bigcup_{\omega \in \Omega} \omega'\) , then \(\mathbf{B}'\) is a basis of a supplementary subspace for \(\mathbf{U}^{\mathrm{H}}\) in \(\mathbf{U}\) .
\end{enumerate}

The union of the set of all \(u_{\omega}\) (for \(\omega \in \Omega\) ) and of \(B'\) is a basis of \(U\) . If \(U' = \sum_{\omega \in \Omega} Au_{\omega}\) and \(U'' = \sum_{b \in B'} Ab\) , we therefore have \(U = U' \oplus U''\) . On the other

hand, we have \(u_{\omega} \in \mathbf{U}^{\mathrm{H}}\) for all \(\omega \in \Omega\) , hence \(\mathbf{U}' \subset \mathbf{U}^{\mathrm{H}}\) . Finally, let \((\alpha_b)_b \in B\) be a family of elements of \(A\) with finite support and let \(x = \sum_{b \in B} \alpha_b b\) . If \(x \in \mathbf{U}^{\mathrm{H}}\) , then

\(\alpha_{hb} = \alpha_b\) for all \(b \in \mathbf{B}\) and all \(h \in \mathbf{H}\) , hence \(x \in \mathbf{U}'\) , and it follows that \(\mathbf{U}' = \mathbf{U}^{\mathrm{H}}\) .

PROPOSITION 5. --- Let M be a free A-module, k an integer \(\geqslant0\) , P the sub-A-module of \(\mathbf{TS}^{k}(\mathbf{M})\) generated by \(\gamma_{k}(\mathbf{M})\) . Assume A to be an infinite integral domain. Then for each \(z\in\mathbf{TS}^{k}(\mathbf{M})\) there exists \(\alpha\in\mathbf{A}-\{0\}\) such that \(\alpha z\in P\) .

Let K be the field of fractions of A. We identify \(\mathbf{T}\mathbf{S}^{k}(\mathbf{M})\) with a sub-A-module of the vector K-space \(\mathbf{V} = \mathbf{T}\mathbf{S}^{k}(\mathbf{M}) \otimes_{\mathbf{A}} \mathbf{K}\) (Prop. 4, and II, p.~314). We have to show that this vector K-space is generated by \(\gamma_{k}(\mathbf{M})\) , that is, every K-linear form f on V satisfying \(f(\gamma_k(M)) = 0\) is zero. Let \((e_i)_{i \in I}\) be a basis of M, and define the \(e_\nu\) as in Prop. 4. For any \((\alpha_i) \in A^{(I)}\) we have, on taking (6) into account,

\[
0 = f \left(\gamma_ {k} \left(\sum_ {i \in I} \alpha_ {i} e _ {i}\right)\right) = \sum_ {\nu \in \mathbf {N} ^ {(I)}, | \nu | = k} \alpha^ {\nu} f (e _ {\nu}).
\]

By Cor. 2 of IV, p.~18 it follows that \(f(e_{\nu}) = 0\) for all \(\nu \in \mathbf{N}^{(I)}\) , whence \(f = 0\) .

